\label{sec:introduccion}

La hipótesis de Riemann relaciona la distribución de los números primos con
la posición de los ceros de una función de variable compleja. Su versión
para funciones $L$ de Dirichlet permite estudiar, además, cómo se reparten
los primos entre las clases de congruencia. Esta introducción construye los
objetos que intervienen en esa relación y prepara la lectura de las trece
notas del cuaderno.

Se presupone familiaridad con series, integrales y álgebra lineal elemental.
Las nociones de análisis complejo y de teoría de números que se utilizan
por primera vez se explican en el texto. Las demostraciones breves se
incluyen; los teoremas clásicos cuya prueba requiere otro desarrollo se
señalan como resultados citados. Estos fundamentos tienen un alcance
propio, independiente de la propuesta que estudia el cuaderno.

\subsection{Primos y funciones de conteo}

Un entero $p\ge2$ es \emph{primo} si sus únicos divisores positivos son
$1$ y $p$. El teorema fundamental de la aritmética afirma que todo entero
$n\ge2$ admite una factorización única en primos, salvo el orden de los
factores. Por ejemplo, $360=2^3 3^2 5$. Esta unicidad será la razón por la
que una serie sobre todos los enteros puede escribirse como un producto
sobre primos.

\begin{proposicion}
Existen infinitos números primos.
\end{proposicion}
\begin{proof}
Si todos ellos fueran $p_1,\ldots,p_r$, el entero
$N=p_1\cdots p_r+1>1$ tendría un divisor primo. Ninguno de los $p_j$
puede dividir $N$, pues entonces dividiría también la diferencia $1$.
Se obtiene una contradicción. El argumento solo necesita un divisor primo
de $N$; no afirma que $N$ sea primo.
\end{proof}

Para medir su distribución definimos, para $x\ge2$,
\[
 \pi(x)=\#\{p\le x:p\text{ primo}\},\qquad
 \vartheta(x)=\sum_{p\le x}\log p.
\]
La función de von Mangoldt reúne los primos y sus potencias:
\[
 \Lambda(n)=
 \begin{cases}
  \log p,&n=p^k,\quad p\text{ primo},\quad k\ge1,\\
  0,&\text{en otro caso}.
 \end{cases}
 \qquad
 \Psi(x)=\sum_{n\le x}\Lambda(n).
\]
Así, $\Lambda(1)=0$, $\Lambda(8)=\log2$ y $\Lambda(12)=0$.
Usamos $\Psi$ mayúscula para esta función de Chebyshev; la letra $\psi$
se reserva para un carácter en las notas posteriores.

El teorema de los números primos, que utilizaremos como resultado clásico,
afirma
\begin{equation}
 \pi(x)\sim\frac{x}{\log x}\qquad(x\longrightarrow\infty).
 \label{eq:intro-pnt}
\end{equation}
La notación $f(x)\sim g(x)$ significa que $f(x)/g(x)\to1$.
Una formulación equivalente de este teorema es $\Psi(x)\sim x$
\cite[§25.16(i)]{dlmf-zeta}. La hipótesis de Riemann trata del tamaño
del error de esa aproximación, que precisaremos después de definir sus
ceros.

Escribiremos $f(x)=O(g(x))$, con $g(x)>0$ para $x$ suficientemente grande,
si existen $C>0$ y $x_0$ tales que $|f(x)|\le Cg(x)$ para todo $x\ge x_0$.
En $O_{q,\varepsilon}$, las constantes pueden depender de los parámetros
indicados. La notación $f(x)=o(g(x))$ expresa $f(x)/g(x)\to0$.
Estas dependencias importan: una estimación para cada $q$ fijo no proporciona
por sí sola una estimación uniforme al variar $q$.

\subsection{El vocabulario de análisis complejo}

Escribimos $s=\sigma+it$, con $\sigma,t\in\R$ e $i^2=-1$.
Su parte real es $\Ree s=\sigma$, su conjugado es $\overline s=\sigma-it$
y su módulo es $|s|=(\sigma^2+t^2)^{1/2}$. Para $x>0$ fijamos
\[
 x^s=e^{s\log x},\qquad |x^s|=x^\sigma.
\]
El logaritmo de $x$ es aquí el logaritmo real; esta convención evita
cualquier elección de rama al escribir $n^{-s}$.

Una función es \emph{holomorfa} en un abierto si admite derivada compleja
en cada punto. La derivada se define mediante el límite del cociente
incremental al aproximarse desde cualquier dirección del plano. Un
resultado básico del análisis complejo garantiza entonces una serie de
Taylor local. Una función \emph{entera} es holomorfa en todo $\C$.

Si una función holomorfa no es idénticamente nula cerca de $\rho$ y se anula
allí, puede escribirse
\[
 f(s)=(s-\rho)^m u(s),\qquad m\ge1,\quad u(\rho)\ne0,
\]
con $u$ holomorfa. El entero $m$ es la \emph{multiplicidad} del cero.
Un \emph{polo de orden $r$} tiene la forma
$f(s)=(s-\rho)^{-r}u(s)$ con $r\ge1$ y $u(\rho)\ne0$.
Una función \emph{meromorfa} es holomorfa salvo polos aislados.

El \emph{residuo} en $\rho$ es el coeficiente de $(s-\rho)^{-1}$ en la
serie de Laurent. Por ejemplo, $1/(s-1)$ tiene un polo simple en $1$ y
residuo $1$. Una singularidad \emph{removible} admite un valor que hace
holomorfa a la función; el cociente $(e^s-1)/s$ se prolonga en $s=0$
con valor $1$.

Una \emph{continuación analítica} prolonga una función a un dominio mayor
conservando sus valores donde ya estaba definida. Su unicidad procede del
teorema de identidad: dos funciones holomorfas en un dominio conexo que
coinciden en un conjunto con un punto de acumulación interior coinciden
en todo el dominio. La expresión que continúa una serie puede existir en
puntos donde la serie original diverge. Usaremos también el hecho clásico
de que un límite uniforme en compactos de funciones holomorfas es
holomorfo \cite[§§1.9--1.10]{dlmf-analisis}.

La fórmula de Cauchy proporciona el control de derivadas que necesitará
la Nota~\ref{sec:ventanas}. Sean $R>0$ y $k\in\Z_{\ge0}$.
Si $f$ es holomorfa en un entorno del disco
cerrado $|z-z_0|\le R$ y $M_R=\max_{|z-z_0|=R}|f(z)|$, entonces
\[
 \frac{f^{(k)}(z_0)}{k!}
   =\frac1{2\pi i}\int_{|z-z_0|=R}\frac{f(z)}{(z-z_0)^{k+1}}\dd z,
 \qquad
 \frac{|f^{(k)}(z_0)|}{k!}\le\frac{M_R}{R^k}.
\]
La integral recorre la circunferencia en sentido antihorario. La fórmula
se cita de \cite[ecuación 1.9.31]{dlmf-analisis}; la cota se deduce
tomando valores absolutos y usando que la circunferencia mide $2\pi R$.
La distancia a una singularidad limita el radio en el que puede aplicarse.

\subsection{La función zeta y el producto de Euler}

Para $\sigma>1$ se define
\begin{equation}
 \zeta(s)=\sum_{n=1}^{\infty}\frac1{n^s}.
 \label{eq:intro-zeta}
\end{equation}
La serie converge absolutamente porque $|n^{-s}|=n^{-\sigma}$.
En un compacto del semiplano $\sigma>1$ existe $\delta>0$ con
$\sigma\ge1+\delta$, y la serie queda dominada por
$\sum n^{-1-\delta}$. La convergencia es uniforme en compactos y
$\zeta$ es holomorfa en ese semiplano.

\begin{proposicion}[Producto de Euler]
Para $\Ree s>1$,
\begin{equation}
 \zeta(s)=\prod_{p\text{ primo}}(1-p^{-s})^{-1},
 \qquad \zeta(s)\ne0.
 \label{eq:intro-euler}
\end{equation}
\end{proposicion}
\begin{proof}
Para un conjunto finito de primos, multiplicar las series geométricas
$1+p^{-s}+p^{-2s}+\cdots$ produce una suma sobre los enteros cuyos
factores primos pertenecen a ese conjunto. La unicidad de la factorización
hace que cada entero aparezca una sola vez. Si incluimos todos los primos
$p\le P$, los términos omitidos corresponden a enteros $n>P$; el valor
absoluto de su suma no supera $\sum_{n>P}n^{-\sigma}$, que tiende a cero.
Esto demuestra el producto.

Para justificar que el producto infinito no se anula, consideramos
\[
 B(s)=\sum_p\sum_{k\ge1}\frac{p^{-ks}}{k}.
\]
La serie converge absolutamente, pues
\[
 \sum_p\sum_{k\ge1}\frac{p^{-k\sigma}}{k}
 \le\frac{1}{1-2^{-\sigma}}\sum_{n\ge2}n^{-\sigma}<\infty.
\]
La identidad de la serie $-\log(1-z)$ para $|z|<1$, aplicada primero a
productos finitos y después al límite, da $\zeta(s)=e^{B(s)}$.
Una exponencial compleja nunca se anula.
\end{proof}

El producto permite recuperar la estructura multiplicativa de los enteros
a partir de una función analítica. Su dominio de convergencia es parte del
enunciado: esta demostración de no anulación se limita a $\sigma>1$.

\subsection{Continuación, ecuación funcional y RH}

Puede verse una primera continuación de $\zeta$ sin desarrollar todavía
su ecuación funcional. Escribamos $\lfloor x\rfloor$ para la parte entera
y $\{x\}=x-\lfloor x\rfloor$ para la parte fraccionaria. Para $\sigma>1$,
la integración término a término da
\[
 \zeta(s)=s\int_1^\infty\lfloor x\rfloor x^{-s-1}\dd x
 =\frac{s}{s-1}-s\int_1^\infty\{x\}x^{-s-1}\dd x.
\]
En efecto, $\lfloor x\rfloor=\sum_{n\ge1}\mathbf1_{[n,\infty)}(x)$,
y $s\int_n^\infty x^{-s-1}\dd x=n^{-s}$; la convergencia absoluta
justifica el intercambio. Como $0\le\{x\}<1$, la última integral converge
y define una función holomorfa en $\sigma>0$. Para derivarla en compactos
se domina cada derivada por una constante por
$(\log x)^k x^{-1-\delta}$, que es integrable. La fórmula prolonga
$\zeta$ a ese semiplano y muestra un polo simple en $s=1$, de residuo $1$.

La expansión local clásica es
$\zeta(s)=(s-1)^{-1}+\gamma_E+O(s-1)$, donde
\[
 \gamma_E=\lim_{N\to\infty}\left(\sum_{n=1}^N\frac1n-\log N\right)
\]
es la constante de Euler. Es la constante denotada por $\gamma$ en la
Nota~\ref{sec:sumas}; no es la función $\Gamma$. La expansión se cita
de \cite[ecuación 25.2.4]{dlmf-zeta}.

La continuación global y la ecuación funcional se utilizan mediante
los siguientes resultados clásicos
\cite[§§25.2, 25.4 y 25.10]{dlmf-zeta}.

\begin{teorema}[Propiedades clásicas de zeta; resultado citado]
La función $\zeta$ admite continuación meromorfa a $\C$, con un único
polo, simple, en $s=1$. La función
\begin{equation}
 \xi(s)=\frac12 s(s-1)\pi^{-s/2}\Gamma(s/2)\zeta(s)
 \label{eq:intro-xi}
\end{equation}
se prolonga a una función entera y satisface $\xi(s)=\xi(1-s)$.
Los ceros de $\zeta$ fuera de $0<\Ree s<1$ son los enteros negativos
pares $-2,-4,-6,\ldots$. Los restantes ceros son los ceros de $\xi$.
\end{teorema}

El factor gamma se define inicialmente por
$\Gamma(z)=\int_0^\infty t^{z-1}e^{-t}\dd t$ para $\Ree z>0$.
La integración por partes da $\Gamma(z+1)=z\Gamma(z)$, y en particular
$\Gamma(n+1)=n!$. Su continuación es meromorfa, tiene polos en los enteros
no positivos y no tiene ceros \cite[§5.2]{dlmf-analisis}. En
\eqref{eq:intro-xi}, los ceros triviales de $\zeta$ cancelan los polos
de $\Gamma(s/2)$ en los enteros negativos pares. El factor $s(s-1)$
cancela los polos restantes en $0$ y $1$; los valores prolongados son
$\xi(0)=\xi(1)=1/2$. La no anulación sobre $\Ree s=1$, utilizada
al delimitar la franja abierta, es también un teorema clásico; no se deduce
solo de la convergencia del producto de Euler.

Los ceros $-2,-4,\ldots$ se llaman \emph{triviales}. Un cero de la franja
$0<\sigma<1$ se llama \emph{no trivial}; $\sigma=1/2$ es la
\emph{recta crítica}. El enunciado de RH es
\begin{equation}
 \boxed{\quad \zeta(\rho)=0,\quad 0<\Ree\rho<1
       \quad\Longrightarrow\quad \Ree\rho=\tfrac12.\quad}
 \label{eq:intro-rh}
\end{equation}

\begin{center}
\setlength{\unitlength}{1mm}
\begin{picture}(135,61)
 \put(30,7){\color{azulnota!9}\rule{60mm}{48mm}}
 \put(8,26){\vector(1,0){118}}
 \put(30,4){\vector(0,1){54}}
 \put(90,7){\color{grisnota}\line(0,1){48}}
 \put(60,7){\color{azulnota}\linethickness{0.8pt}\line(0,1){48}}
 \put(125,21){\makebox(0,0)[l]{$\sigma$}}
 \put(26,58){\makebox(0,0)[r]{$t$}}
 \put(28,21){\makebox(0,0)[r]{$0$}}
 \put(58,21){\makebox(0,0)[r]{$1/2$}}
 \put(88,21){\makebox(0,0)[r]{$1$}}
 \put(90,26){\circle*{1.6}}
 \put(94,30){\makebox(0,0)[l]{\small polo de $\zeta$}}
 \put(60,59){\makebox(0,0){\small franja crítica}}
 \put(62,47){\makebox(0,0)[l]{\small recta crítica}}
 \put(60,1){\makebox(0,0){\small La franja se extiende indefinidamente en ambas direcciones verticales.}}
\end{picture}
\end{center}

La conjugación y la ecuación funcional producen las simetrías
$\rho\mapsto\overline\rho$ y $\rho\mapsto1-\rho$. Una figura simétrica
puede contener puntos a ambos lados de su eje: la simetría por sí sola no
demuestra RH. Lo que se debe excluir es cualquier cero no trivial con
parte real distinta de $1/2$.

\subsection{Congruencias y caracteres de Dirichlet}

La congruencia $a\equiv b\pmod q$ significa $q\mid(a-b)$. Para $q\ge2$,
las clases coprimas con $q$ forman el grupo multiplicativo
\[
 U_q=(\Z/q\Z)^\times,\qquad |U_q|=\varphi(q).
\]
La función $\varphi$ de Euler cuenta esas clases. Por ejemplo,
$U_4=\{1,3\}$ y $\varphi(4)=2$. Un grupo tiene una operación asociativa,
un elemento neutro e inversos; es abeliano cuando su operación es
conmutativa. Un homomorfismo conserva la operación y el elemento neutro.
Un \emph{carácter} de un grupo abeliano
finito $G$ es un homomorfismo $G\to\C^\times$. Sus valores son raíces
de la unidad: si $g^r=1$, entonces $\chi(g)^r=1$.

\begin{definicion}
Un carácter de Dirichlet módulo $q$ es un carácter de $U_q$, extendido
a los enteros mediante $\chi(n)=0$ cuando $\gcd(n,q)>1$ y mediante
su valor en la clase de $n$ cuando $\gcd(n,q)=1$. Para $q=1$ adoptamos
el grupo trivial, $\varphi(1)=1$ y el único carácter $\chi(n)=1$.
\end{definicion}

La extensión es periódica y completamente multiplicativa:
$\chi(mn)=\chi(m)\chi(n)$. Su valor absoluto es $1$ en los enteros coprimos con
$q$ y $0$ en los demás. El \emph{carácter principal} $\chi_{0,q}$ vale
$1$ en las unidades y $0$ fuera de ellas. No se debe identificar
$\chi_{0,q}$, para $q>1$, con la función constante $1$.

Los dos caracteres módulo $4$ son
\[
\begin{array}{c|rrrr}
 n\bmod4&0&1&2&3\\ \hline
 \chi_{0,4}(n)&0&1&0&1\\
 \chi_4(n)&0&1&0&-1
\end{array}
\]
El segundo distingue las dos clases impares. Su cuadrado es el carácter
principal módulo $4$. El \emph{orden} de un carácter es su orden en el
grupo de caracteres: el menor $o\ge1$ con $\chi^o=\chi_{0,q}$.
La potencia de exponente cero se entiende también como ese carácter
principal. Estas definiciones y la ortogonalidad se recogen en
\cite[§27.8]{dlmf-dirichlet}.

\subsection{Conductor, carácter primitivo y función \texorpdfstring{$L$}{L}}

Si $f\mid q$, un carácter módulo $q$ es \emph{inducido} por uno módulo
$f$ cuando sus valores en las unidades módulo $q$ se obtienen reduciendo
módulo $f$. Su \emph{conductor} es el menor módulo divisor de $q$ desde
el que se induce. Cada carácter procede de un único carácter
\emph{primitivo} $\chi^*$ de ese conductor; un carácter es primitivo cuando
su módulo coincide con su conductor. La existencia y unicidad de este
carácter primitivo se citan de \cite[Teo.~18.13 y Def.~18.14]{sutherland-dirichlet}.

El carácter $\chi_4$ es primitivo de conductor $4$. Si se lo considera
como carácter inducido módulo $8$, conserva los mismos valores en todos
los enteros, pero su módulo es $8$ y su conductor sigue siendo $4$.
El carácter principal de cualquier módulo tiene conductor $1$.
En particular, al elevar un carácter primitivo a una potencia puede
disminuir el conductor. En las notas, $\psi^{[j]}$ designa el carácter
primitivo que induce a $\psi^j$.

Para $\Ree s>1$ se define la función $L$ de Dirichlet por
\begin{equation}
 L(s,\chi)=\sum_{n\ge1}\frac{\chi(n)}{n^s}
          =\prod_p(1-\chi(p)p^{-s})^{-1}.
 \label{eq:intro-l}
\end{equation}
La convergencia absoluta se deduce de $|\chi(n)|\le1$. La prueba del
producto de Euler se repite usando la multiplicatividad de $\chi$;
la serie de logaritmos converge absolutamente, y por ello
$L(s,\chi)\ne0$ en $\Ree s>1$.

Al inducir desde un conductor menor se suprimen factores de Euler:
\begin{equation}
 L(s,\chi)=L(s,\chi^*)
   \prod_{\substack{p\mid q\\p\nmid f}}(1-\chi^*(p)p^{-s}).
 \label{eq:intro-inducido}
\end{equation}
Esta igualdad se demuestra en $\Ree s>1$ comparando los productos, y se
extiende por continuación meromorfa. Para el carácter principal,
\[
 L(s,\chi_{0,q})=\zeta(s)\prod_{p\mid q}(1-p^{-s}).
\]
Su residuo en $1$ es $\prod_{p\mid q}(1-p^{-1})=\varphi(q)/q$.
Para $q=1$ el producto es vacío y se recupera $\zeta$.

Inducir $\chi_4$ a módulo $8$ no elimina ningún factor adicional.
Inducirlo a módulo $12$ elimina el factor de $3$ y da
$L(s,\chi)=L(s,\chi_4)(1+3^{-s})$. Los posibles ceros de los factores
añadidos en \eqref{eq:intro-inducido} están sobre $\Ree s=0$, porque
$|\chi^*(p)|=1$. Por eso los ceros de la franja abierta se estudian
con los caracteres primitivos.

\begin{teorema}[Propiedades clásicas de las funciones $L$; resultado citado]
Si $\chi$ no es principal, $L(s,\chi)$ se prolonga a una función entera
y $L(1,\chi)\ne0$. Para todo carácter, $L(s,\chi)$ no se anula en la
recta $\Ree s=1$; en el caso principal, $s=1$ es un polo.
\end{teorema}

La continuación y el valor no nulo en $1$ se citan de
\cite[§25.15, en particular (25.15.9)]{dlmf-dirichlet}; la no anulación
en el resto de la recta, de \cite[§§3.3--3.4]{kedlaya-ant}. El valor no nulo
en $1$ permite definir $L'(1,\chi)/L(1,\chi)$ en la Nota~\ref{sec:sumas}.
En el carácter primitivo principal aparece el polo de $\zeta$, que exige
una corrección diferente.

\subsection{Ecuación funcional y enunciado de GRH}

Sea $\chi$ primitivo de conductor $q>1$. Su \emph{paridad} es
$a_\chi\in\{0,1\}$, determinada por $\chi(-1)=(-1)^{a_\chi}$.
Definimos la función completada
\begin{equation}
 \mathcal L(s,\chi)=
 \left(\frac q\pi\right)^{(s+a_\chi)/2}
 \Gamma\!\left(\frac{s+a_\chi}{2}\right)L(s,\chi).
 \label{eq:intro-completada}
\end{equation}
La letra $\mathcal L$ evita confundir esta función con $\Lambda(n)$.
El teorema clásico de la ecuación funcional afirma que $\mathcal L$ es
entera y satisface
\begin{equation}
 \mathcal L(s,\chi)=\varepsilon_\chi\mathcal L(1-s,\overline\chi),
 \qquad
 \varepsilon_\chi=\frac{\tau(\chi)}{i^{a_\chi}\sqrt q},
 \quad |\varepsilon_\chi|=1,
 \label{eq:intro-fe-l}
\end{equation}
donde $\tau(\chi)=\sum_{r=1}^{q}\chi(r)e^{2\pi ir/q}$ es la suma de
Gauss. La primitividad es una hipótesis de esta formulación. Citamos la
ecuación funcional y las propiedades de sus ceros de
\cite[§25.15]{dlmf-dirichlet}.

Los polos del factor gamma fuerzan ceros triviales de $L$ en
$0,-2,-4,\ldots$ si el carácter es par y no principal, y en
$-1,-3,-5,\ldots$ si es impar. Los demás ceros están en
$0<\Ree s<1$. El caso de conductor $1$ es $\zeta$ y utiliza
$\xi$ de \eqref{eq:intro-xi}; sus ceros triviales comienzan en $-2$.

\begin{definicion}[GRH para funciones $L$ de Dirichlet]
Para todo carácter primitivo $\chi$ de conductor $q\ge1$ y todo
$\rho\in\C$,
\begin{equation}
 L(\rho,\chi)=0,\qquad 0<\Ree\rho<1
 \quad\Longrightarrow\quad \Ree\rho=\frac12.
 \label{eq:intro-grh}
\end{equation}
\end{definicion}

Esta es la acepción de GRH utilizada en el cuaderno. Incluye RH mediante
el carácter de conductor $1$. Las funciones $L$ automorfas más generales
requieren otros enunciados.

\begin{proposicion}[Reducción al lado derecho de la franja]
Para probar \eqref{eq:intro-grh} basta excluir, para todos los caracteres
primitivos, los ceros con $1/2<\Ree\rho<1$.
\end{proposicion}
\begin{proof}
Si $L(\rho,\chi)=0$ y $0<\Ree\rho<1/2$, los factores que completan
$L$ son finitos y no nulos en $\rho$. La ecuación funcional produce un
cero en $1-\rho$ de la función asociada a $\overline\chi$, que también
es primitivo. Su parte real es mayor que $1/2$. Para conductor $1$ se
aplica $\xi(s)=\xi(1-s)$. Esto contradice la exclusión supuesta.
\end{proof}

Para un carácter complejo, la conjugación relaciona dos funciones:
$L(\overline s,\overline\chi)=\overline{L(s,\chi)}$.
Por eso se conserva $\overline\chi$ en la ecuación funcional. No se
presupone simetría respecto del eje real para los ceros de un carácter
complejo fijo.

\subsection{Ortogonalidad y primos en progresiones}

Los caracteres permiten extraer una clase de congruencia mediante una
identidad finita. El conjunto de caracteres de $U_q$ se denota
$\widehat U_q$ y tiene $\varphi(q)$ elementos. El resultado algebraico
que asegura que los caracteres separan los puntos se demuestra en la
Nota~\ref{sec:caracteres}.

\begin{proposicion}[Ortogonalidad]
Si $\gcd(a,q)=1$, entonces, para todo entero $n$,
\begin{equation}
 \mathbf1_{n\equiv a\ (\mathrm{mod}\ q)}
  =\frac1{\varphi(q)}\sum_{\chi\in\widehat U_q}
       \overline{\chi(a)}\chi(n).
 \label{eq:intro-ortogonalidad}
\end{equation}
\end{proposicion}
\begin{proof}
Si $n$ no es coprimo con $q$, ambos miembros son cero. Si lo es, ponemos
$g=na^{-1}\in U_q$. Para $g=1$, todos los sumandos valen $1$.
Para $g\ne1$ existe un carácter $\lambda$ con $\lambda(g)\ne1$.
La multiplicación por $\lambda$ permuta $\widehat U_q$, de modo que
\[
 \sum_\chi\chi(g)=\sum_\chi(\lambda\chi)(g)
 =\lambda(g)\sum_\chi\chi(g).
\]
La suma es cero. Esto demuestra los dos casos.
\end{proof}

Definiendo
\[
 \Psi(x,\chi)=\sum_{n\le x}\Lambda(n)\chi(n),\qquad
 \Psi(x;q,a)=\sum_{\substack{n\le x\\n\equiv a\ (\mathrm{mod}\ q)}}\Lambda(n),
\]
la identidad anterior implica
\begin{equation}
 \Psi(x;q,a)=\frac1{\varphi(q)}
  \sum_{\chi\in\widehat U_q}\overline{\chi(a)}\Psi(x,\chi).
 \label{eq:intro-progresiones}
\end{equation}
Por ejemplo, para $q=4$, la diferencia entre los conteos ponderados en
las clases $1$ y $3$ es $\Psi(x,\chi_4)$.

El teorema de Dirichlet afirma que cada progresión $a+kq$ con
$\gcd(a,q)=1$ contiene infinitos primos
\cite[Teo.~18.1]{sutherland-dirichlet}. El teorema de los números primos
en progresiones precisa, para $q$ fijo,
\[
 \pi(x;q,a):=\#\{p\le x:p\equiv a\pmod q\}
 \sim\frac1{\varphi(q)}\frac{x}{\log x},
 \qquad
 \Psi(x;q,a)\sim\frac{x}{\varphi(q)}.
\]
Son resultados clásicos citados; la identidad finita
\eqref{eq:intro-progresiones} explica el papel de los caracteres, pero
la estimación de las sumas exige análisis adicional
\cite[Teo.~4.12]{kedlaya-ant}.

La relación clásica entre RH y el error de conteo puede expresarse como
\begin{equation}
 \mathrm{RH}\quad\Longleftrightarrow\quad
 \text{para todo }\varepsilon>0,\qquad
 \Psi(x)=x+O_\varepsilon(x^{1/2+\varepsilon}).
 \label{eq:intro-error-rh}
\end{equation}
Citamos esta equivalencia de \cite[§25.16(i)]{dlmf-zeta}.
Bajo GRH, para cada $q$ fijo y cada $a$ coprimo con $q$, se obtiene
\[
 \Psi(x;q,a)=\frac{x}{\varphi(q)}
     +O_{q,\varepsilon}(x^{1/2+\varepsilon})
 \qquad\text{para todo }\varepsilon>0.
\]
Esta consecuencia se obtiene de la cota clásica
$\Psi(x,\chi)-\delta_\chi x=O(\sqrt{x}\log^2(qx))$ bajo GRH, donde
$\delta_\chi=1$ para el carácter principal y $0$ para los demás
\cite[pp.~19--20 del PDF]{vaughan-grh}. Se aplica
\eqref{eq:intro-progresiones} y, para $q$ fijo, se absorbe
$\log^2(qx)$ en $O_{q,\varepsilon}(x^\varepsilon)$.
El subíndice conserva la
dependencia permitida en $q$. Tampoco se ha sustituido el exponente
$1/2+\varepsilon$ por una afirmación sin margen.

Si $\rho=\beta+i\gamma$, entonces
$x^\rho=x^\beta e^{i\gamma\log x}$. El primer factor determina una
escala de tamaño y el segundo una oscilación. Las fórmulas explícitas de
la teoría relacionan términos de esta forma con los ceros; esta observación
explica el interés de la escala $x^{1/2}$. Por sí sola no justifica sumar
una serie sobre todos los ceros ni estimar sus cancelaciones.

\subsection{De los ceros a los polos: derivada logarítmica y Mellin}

El producto de Euler lleva a una identidad que reaparece en todo el
cuaderno:
\begin{equation}
 -\frac{L'(s,\chi)}{L(s,\chi)}
   =\sum_{n\ge1}\frac{\Lambda(n)\chi(n)}{n^s},
 \qquad \Ree s>1.
 \label{eq:intro-logderivada}
\end{equation}
Para derivarla, se utiliza la serie convergente
$\log L(s,\chi)=\sum_p\sum_{k\ge1}\chi(p)^k/(k p^{ks})$,
entendida como un logaritmo holomorfo definido por esa serie.
La derivada de $p^{-ks}$ aporta $-k\log p$ y cancela el denominador $k$.
Quedan los términos $(\log p)\chi(p)^k p^{-ks}$ correspondientes a
$n=p^k$. La convergencia uniforme de la serie derivada en compactos de
$\Ree s>1$ se controla con $\sum (\log n)n^{-1-\delta}<\infty$.

\begin{proposicion}
Si $\rho$ es un cero de multiplicidad $m$ de $L(s,\chi)$, la función
$-L'/L$ tiene en $\rho$ un polo simple de residuo $-m$.
\end{proposicion}
\begin{proof}
Escribimos $L(s,\chi)=(s-\rho)^m u(s)$, con $u(\rho)\ne0$.
Entonces
\[
 -\frac{L'(s,\chi)}{L(s,\chi)}
 =-\frac{m}{s-\rho}-\frac{u'(s)}{u(s)},
\]
y el segundo sumando es holomorfo cerca de $\rho$.
\end{proof}

La transformada de Mellin convierte escalas multiplicativas en información
analítica. En la convención usada aquí se considera
\[
 \mathcal M_1 f(s)=\int_1^\infty f(x)x^{-s-1}\dd x.
\]
El cambio $x=e^t$ la convierte en
$\int_0^\infty f(e^t)e^{-st}\dd t$, una transformada de Laplace.
Su dominio inicial depende del crecimiento de $f$.
La convención habitual $\int_0^\infty f(x)x^{z-1}\dd x$ se recupera
extendiendo $f$ por cero por debajo de $1$ y poniendo $z=-s$
\cite[ecuación 2.5.1]{dlmf-analisis}.

La Nota~\ref{sec:sumas} demuestra, con intercambio justificado de suma
e integral, que la suma suavizada
\[
 S_\chi(x)=\sum_{n\le x}\Lambda(n)\chi(n)(x/n-1)
\]
satisface
\[
 \mathcal M_1 S_\chi(s)=-\frac{L'(s,\chi)}{s(s-1)L(s,\chi)}
 \qquad(\Ree s>1).
\]
El peso se anula cuando $x=n$ y suaviza la entrada de cada término.
Tras corregir la singularidad en $1$, un cero no trivial $\rho$ sigue
produciendo un polo, cuyo residuo es
$m/[\rho(1-\rho)]$. Este factor suele ser complejo; la positividad
que se utilizará es la de las multiplicidades y de los coeficientes
que lo multiplican, no una afirmación de signo sobre cada residuo.

Para un carácter primitivo $\psi$ de orden $o$, se utiliza la combinación
\[
 H_\psi(x)=G_1(x)+\frac1{3o+1}
     \sum_{j=1}^{o-1}\Ree G_{\psi^{[j]}}(x),
\]
en la que $G_\chi$ es la suma corregida definida en la Nota~\ref{sec:sumas}.
El objetivo aritmético de las fuentes exige
\begin{equation}
 \forall\eta>0\ \forall A\ge2\ \exists Y\ge A\ \forall x\in[Y,Y^{3/2}],
 \qquad H_\psi(x)\le x^{1/2+\eta}.
 \label{eq:intro-ventanas}
\end{equation}
La Nota~\ref{sec:ventanas} desarrolla el criterio analítico que utiliza
esas ventanas para excluir polos. Una cota superior unilateral en algunos
intervalos y una estimación global en valor absoluto, como la que aparece
en \eqref{eq:intro-error-rh}, tienen cuantificadores distintos.
La producción aritmética de las ventanas pertenece a los pasos de las
fuentes que siguen por estudiar.

\subsection{Retículos, índices y pendientes: una entrada lineal}

Un \emph{retículo pleno} en un espacio vectorial racional $V$ de dimensión
$r$ es un subgrupo $M$ generado por una base racional:
$M=\Z v_1\oplus\cdots\oplus\Z v_r$. Su rango es $r$ y
$M\otimes_\Z\Q=V$. La estructura integral retiene información que se
pierde al permitir todos los coeficientes racionales.

Si $M'\subset M$ tiene el mismo rango, el \emph{índice}
$[M:M']$ es el número finito de elementos de $M/M'$. Por ejemplo,
$\Z\times p\Z\subset\Z^2$ tiene índice $p$. Un operador puede preservar
$M$ mientras que sus proyectores espectrales solo estén definidos sobre
$\Q$: la Nota~\ref{sec:proyecciones} muestra el operador
$T=\left(\begin{smallmatrix}0&1\\0&p\end{smallmatrix}\right)$ y el
proyector $T/p$. El vector $(0,1)$ se proyecta en $(1/p,1)$.

Una norma euclídea en $V\otimes_\Q\R$ permite medir el volumen del
paralelepípedo generado por una base de $M$, llamado \emph{covolumen}.
Un cambio de base integral tiene determinante $\pm1$, por lo que ese
volumen no depende de la base elegida.
Como modelo sobre $\Q$, para $r>0$ ponemos
\[
 \widehat{\deg}(M)=-\log\operatorname{covol}(M),\qquad
 \widehat\mu(M)=\frac{\widehat{\deg}(M)}r.
\]
La segunda cantidad es la \emph{pendiente media}. Si $M'\subset M$
tiene índice $d$ y se utiliza la misma métrica, su covolumen se multiplica
por $d$ y
\[
 \widehat\mu(M')=\widehat\mu(M)-\frac{\log d}{r}.
\]
Un índice y un denominador pueden así convertirse en un coste logarítmico.
Las fuentes trabajan con retículos sobre anillos de enteros y normas en
las distintas inmersiones; esta definición euclídea es una entrada a esas
normalizaciones. En particular, A y B usan la pendiente duplicada
$\mu=2\mu_{\mathrm{Ar}}$: sobre $\Q$, con la misma métrica, es el doble
de la pendiente media definida aquí. El factor debe conservarse al
comparar sus desigualdades \cite[§5, (5.1)--(5.3), pp.~17--18]{fuenteA}
\cite[§4.3, p.~39]{fuenteB}.

En el lema local se utiliza un \emph{anillo de valoración discreta}
$\mathcal O$: sus elementos no nulos del cuerpo de fracciones se escriben
$u\varpi^k$, con $u$ invertible en $\mathcal O$ y $k\in\Z$ único.
La \emph{uniformizante} $\varpi$ genera su ideal maximal.
La valoración es $v(u\varpi^k)=k$ y $\mathcal O$ consta de $0$ y los
elementos de valoración no negativa. La \emph{longitud} de un módulo
de longitud finita cuenta los factores simples de una serie de composición.
Los factores simples son copias del cuerpo residual $\mathcal O/(\varpi)$.
Por ejemplo, la cadena de potencias de $\varpi$ muestra que
$\mathcal O/(\varpi^k)$ tiene longitud $k$, para $k\ge1$.
En $\Z_p$, el anillo de enteros $p$-ádicos, cada factor simple es
$\mathbf F_p$; un módulo de longitud $k$ tiene entonces $p^k$ elementos.
La Nota~\ref{sec:proyecciones} desarrolla la cota que necesita este lenguaje.
Las definiciones de valoración y anillo local se encuentran en
\cite[Lecture 1, §§1.3--1.5]{sutherland-algebra}.

\subsection{Cuerpos de números, inercia y Frobenius}

Un \emph{cuerpo de números} es una extensión finita $K/\Q$.
Su grado $[K:\Q]$ es su dimensión como espacio vectorial racional.
El anillo $\mathcal O_K$ de \emph{enteros algebraicos} reúne los elementos
de $K$ que satisfacen un polinomio mónico con coeficientes en $\Z$.
Un ideal de $\mathcal O_K$ es un subgrupo aditivo estable bajo
multiplicación por elementos de $\mathcal O_K$. Un ideal propio
$\mathfrak p$ es \emph{primo} si $ab\in\mathfrak p$ implica
$a\in\mathfrak p$ o $b\in\mathfrak p$. Decimos que está \emph{sobre $p$}
cuando $\mathfrak p\cap\Z=p\Z$, y escribimos $\mathfrak p\mid p$.
Para un ideal primo no nulo $\mathfrak p$, el cociente $\mathcal O_K/\mathfrak p$ es un
cuerpo finito; su cardinal es la norma $N\mathfrak p$.

La factorización de un primo racional se expresa como
\[
 p\mathcal O_K=\prod_{j=1}^g\mathfrak p_j^{e_j},\qquad
 N\mathfrak p_j=p^{f_j},\qquad
 \sum_{j=1}^g e_jf_j=[K:\Q].
\]
El entero $e_j$ es el índice de ramificación y $f_j$ el grado residual.
Estas propiedades de los anillos de enteros se utilizan como resultados
clásicos \cite{milne-ant}. En $\Q(i)$, cuyo anillo de enteros es
$\Z[i]$, el primo $5$ se descompone como $(2+i)(2-i)$, mientras que
$2=-i(1+i)^2$ se ramifica. La factorización es una factorización de
ideales en el enunciado general; no se supone factorización única de
elementos en cualquier $\mathcal O_K$.

Supongamos que $K/\Q$ es de Galois: en característica cero, esto equivale
a que los polinomios irreducibles de $\Q[X]$ con una raíz en $K$
se descompongan allí.
El grupo $G=\operatorname{Gal}(K/\Q)$ consta de los automorfismos de $K$
que fijan $\Q$. Para $\mathfrak p\mid p$, su grupo de descomposición
$D_{\mathfrak p}$ es el estabilizador de $\mathfrak p$. El grupo de
\emph{inercia} $I_{\mathfrak p}$ es el núcleo de la acción de
$D_{\mathfrak p}$ sobre el cuerpo residual.

El cociente $D_{\mathfrak p}/I_{\mathfrak p}$ se identifica con el grupo
de Galois del cuerpo residual sobre $\mathbf F_p$. Su generador
$u\mapsto u^p$ define el \emph{Frobenius aritmético}. Cuando el primo
no se ramifica, la inercia es trivial y este elemento se identifica con
un elemento de $D_{\mathfrak p}$. Cambiar el primo de $K$ sobre $p$
lo conjuga; si $G$ es abeliano y el primo no se ramifica, el elemento
no depende de esa elección \cite[Lecture 7, §7.3]{sutherland-algebra}.

Para un carácter unidimensional $\chi$ de $D_{\mathfrak p}$ y su
subgrupo de inercia $I=I_{\mathfrak p}$,
\[
 \frac1{|I|}\sum_{u\in I}\chi(u)=
 \begin{cases}
  1,&\chi|_I=1,\\
  0,&\chi|_I\ne1.
 \end{cases}
\]
La demostración vuelve a ser la cancelación de una suma de caracteres
por traslación. Este promedio selecciona los caracteres que descienden
a $D_{\mathfrak p}/I_{\mathfrak p}$. En el caso abeliano de la
Nota~\ref{sec:caracteres} también se aplica a caracteres de $G$ y al
cociente $G/I_{\mathfrak p}$. Así se incorpora la ramificación en la
identidad finita de la Nota~\ref{sec:caracteres}.

Los periodos automorfos, los operadores de Hecke y la cohomología que
aparecen después suministran mapas, condiciones de integridad y cotas.
Su construcción necesita un estudio adicional. La filtración de
Harder--Narasimhan organiza subobjetos por pendientes; la desigualdad de
marcos controla cómo una familia de mapas cubre direcciones del destino.
La guía que sigue localiza estos resultados en A y B, y las notas
distinguen sus hipótesis de los pasos ya demostrados.

\subsection{Ejercicios de lectura y soluciones}

\begin{ejercicio}\label{ej:intro-mangoldt}
Calcule $\pi(10)$, $\vartheta(10)$ y $\Psi(10)$ e indique qué términos
explican la diferencia entre las dos últimas funciones.
\end{ejercicio}
\begin{ejercicio}\label{ej:intro-euler}
Para $\Ree s>1$, describa los enteros que aparecen al expandir
$(1-2^{-s})^{-1}(1-3^{-s})^{-1}$. Explique por qué esta expresión aún
no es $\zeta(s)$.
\end{ejercicio}
\begin{ejercicio}\label{ej:intro-conductor}
Induzca $\chi_4$ a módulo $12$. Calcule su valor en $3$ y en $5$, su
conductor y el factor que relaciona su función $L$ con $L(s,\chi_4)$.
\end{ejercicio}
\begin{ejercicio}\label{ej:intro-polo}
Si $f(s)=(s-\rho)^3e^s$, calcule $-f'(s)/f(s)$ y su residuo en $\rho$.
Si $0<\Ree\rho<1$, calcule el residuo después de dividir por $s(s-1)$.
\end{ejercicio}
\begin{ejercicio}\label{ej:intro-ortogonalidad}
Escriba las indicadoras de las clases $1$ y $3$ módulo $4$ usando
$\chi_{0,4}$ y $\chi_4$, y compruebe sus valores en un entero par.
\end{ejercicio}
\begin{ejercicio}\label{ej:intro-reticulo}
Con la métrica euclídea habitual, calcule índice, covolumen, grado y
pendiente media de $M'=\Z\times p\Z\subset\Z^2$.
\end{ejercicio}

\paragraph{Solución del Ejercicio~\ref{ej:intro-mangoldt}.}
Los primos son $2,3,5,7$, por lo que $\pi(10)=4$ y
$\vartheta(10)=\log2+\log3+\log5+\log7$.
Se añaden las potencias $4,8,9$, de modo que
$\Psi(10)=3\log2+2\log3+\log5+\log7$.

\paragraph{Solución del Ejercicio~\ref{ej:intro-euler}.}
La expansión es $\sum_{a,b\ge0}(2^a3^b)^{-s}$. Cada entero con factores
primos únicamente $2$ y $3$ aparece una vez. Faltan, por ejemplo, $5$, $7$
y todos sus múltiplos que contienen factores primos diferentes de $2$ y
$3$. Para obtener zeta se toma el límite de los productos finitos, en
el dominio donde la convergencia absoluta permite hacerlo.

\paragraph{Solución del Ejercicio~\ref{ej:intro-conductor}.}
El carácter inducido vale cero en $3$, porque $\gcd(3,12)>1$, y vale
$\chi_4(5)=1$ en $5$. Su conductor es $4$. Como $\chi_4(3)=-1$,
se obtiene $L(s,\chi)=L(s,\chi_4)(1+3^{-s})$.

\paragraph{Solución del Ejercicio~\ref{ej:intro-polo}.}
Se tiene $-f'/f=-3/(s-\rho)-1$, con residuo $-3$. El factor
$1/[s(s-1)]$ es holomorfo en $\rho$ bajo la hipótesis indicada, así que
el nuevo residuo es $-3/[\rho(\rho-1)]=3/[\rho(1-\rho)]$.

\paragraph{Solución del Ejercicio~\ref{ej:intro-ortogonalidad}.}
Las indicadoras son $(\chi_{0,4}+\chi_4)/2$ y
$(\chi_{0,4}-\chi_4)/2$, respectivamente. En un entero par ambos
caracteres valen cero y las dos indicadoras también. Sustituir
$\chi_{0,4}$ por la función constante $1$ daría un valor incorrecto.

\paragraph{Solución del Ejercicio~\ref{ej:intro-reticulo}.}
El cociente $\Z^2/M'$ es isomorfo a $\Z/p\Z$, luego el índice es $p$.
La base $(1,0),(0,p)$ determina un rectángulo de área $p$; ese es su
covolumen. El grado es $-\log p$ y, como el rango es $2$, la pendiente
media es $-(\log p)/2$.
